\section{Reproduction Setup}\label{app:setup}

\subsection{Oracle AI Database 26ai Free container}

The evaluation database is the official Free container image, started with the demo password used only for this report:

\begin{lstlisting}[style=codestyle]
docker run -d --name oraviz-oracle -p 1530:1521 \
  -e ORACLE_PWD=OraViz2026 \
  container-registry.oracle.com/database/free:latest
\end{lstlisting}

The demo schema is loaded from the SQL script shipped in the repository (\texttt{examples/demo-sales.sql}). It creates the application user, the \texttt{SALES\_DEMO} table with 96 deterministic rows (12 months, 4 regions, 2 channels), and \texttt{PRODUCT\_VECTORS} with an 8-dimensional \texttt{VECTOR} column:

\begin{lstlisting}[style=codestyle]
docker exec -i oraviz-oracle sqlplus -S oraviz/OraViz2026@//localhost:1521/FREEPDB1 \
  < examples/demo-sales.sql
\end{lstlisting}

\subsection{Benchmark and figures}

The benchmark needs the dev extra (which adds \texttt{tiktoken}) and a SQLcl 26.1 installation for the official server:

\begin{lstlisting}[style=codestyle]
uv run --extra dev python benchmarks/run_benchmarks.py \
  --dsn localhost:1530/FREEPDB1 --user oraviz --password OraViz2026 \
  --sqlcl /path/to/sqlcl/bin/sql
\end{lstlisting}

Results are written to \texttt{benchmarks/results/benchmark-results.json} and \texttt{.md}; the figures in this report are regenerated from the JSON with \texttt{uv run python paper/make\_figures.py}.

\subsection{Evaluated versions}

\begin{table}[h]
\centering
\small
\begin{tabular}{ll}
\toprule
Component & Version \\
\midrule
Oracle AI Database 26ai Free & 23.26.0 container image \\
SQLcl MCP server & 26.1.0 \\
OraViz MCP server & 0.1.0 (\texttt{main}) \\
Python & 3.12.3 \\
FastMCP & 4.0.3 \\
python-oracledb & 26.0.0 (thin mode) \\
Matplotlib & 3.11.2 (Agg backend) \\
tiktoken & 0.14.0 (\texttt{cl100k\_base}, \texttt{o200k\_base}) \\
Java (for SQLcl) & 21 \\
\bottomrule
\end{tabular}
\caption{Evaluated versions. Formatting of official-server output can change between SQLcl releases, so the numbers are tied to this configuration.}
\label{tab:versions}
\end{table}
